% This is the ADASS_template.tex LaTeX file, 24 Oct 2025.
% It is based on the ASP general author template file, but modified to reflect the specific
% requirements of the ADASS proceedings.

% This is the aspauthor.tex LaTeX file
% Copyright 2021, Astronomical Society of the Pacific Conference Series
% Revision:  29 November 2021

% To compile, at the command line positioned at this folder, type:
% pdflatex ADASS_template
% pdflatex ADASS_template
% This will create a file called ADASS_template.pdf

\documentclass[11pt,twoside]{article}

% Do NOT use ANY packages other than asp2023. 
\usepackage{asp2023}

\aspSuppressVolSlug
\resetcounters

% References must all use BibTeX entries in a .bibfile.
% References must be cited in the text using \citet{} or \citep{}.
% Do not use \cite{}.
% See ManuscriptInstructions.pdf for more details
\bibliographystyle{asp2021}

% The ``markboth'' line sets up the running heads for the paper.
% 1 author: "Surname"
% 2 authors: "Surname1 and Surname2"
% 3 authors: "Surname1, Surname2, and Surname3"
% >3 authors: "Surname1 et al."
% Replace ``Short Title'' with the actual paper title, shortened if necessary.
% Use mixed case type for the shortened title
% Ensure shortened title does not cause an overfull hbox LaTeX error
% See ASPmanual2010.pdf 2.1.4  and ManuscriptInstructions.pdf for more details
\markboth{Bryant et al.}{ATRAN\_SDC Atmospheric Transmission Model}

\begin{document}

\title{ATRAN\_SDC - Modelling Atmospheric Transmission Spectra for Airborne \& Ground-Based Infrared Observations}

% Note the position of the comma between the author name and the 
% affiliation number.
% Authors surnames should come after first names or initials, eg John Smith, or J. Smith.
% Author names should be separated by commas.
% The final author should be preceded by "and".
% Affiliations should not be repeated across multiple \affil commands. If several
% authors share an affiliation this should be in a single \affil which can then
% be referenced for several author names. If only one affiliation, no footnotes are needed.
% See ManuscriptInstructions.pdf and ASP's manual2010.pdf 3.1.4 for more details
\author{Aaron Bryant, Jonas Fr\"{u}h, Benjamin Greiner, Christof Iserlohe, Bastian Knieling, Mario Ruoff, and Bernhard Schulz}
\affil{SOFIA Data Center, Institut f\"{u}r Raumfahrtsysteme, Universit\"{a}t Stuttgart, Pfaffenwaldring 31, 70569 Stuttgart, Germany; \email{aaron.bryant@irs.uni-stuttgart.de}}

% This section is for ADS Processing.  There must be one line per author. paperauthor has 9 arguments.
\paperauthor{Aaron Bryant}{aaron.bryant@irs.uni-stuttgart.de}{0000-0002-9472-2391}{Universit\"{a}t Stuttgart}{Institut f\"{u}r Raumfahrtsysteme}{Stuttgart}{Baden-W\"{u}rttemberg}{70569}{Germany}
\paperauthor{Jonas Fr\"{u}h}{jonas.frueh@irs.uni-stuttgart.de}{0009-0006-4927-3861}{Universit\"{a}t Stuttgart}{Institut f\"{u}r Raumfahrtsysteme}{Stuttgart}{Baden-W\"{u}rttemberg}{70569}{Germany}
\paperauthor{Benjamin Greiner}{benjamin.greiner@irs.uni-stuttgart.de}{0009-0008-1244-7711}{Universit\"{a}t Stuttgart}{Institut f\"{u}r Raumfahrtsysteme}{Stuttgart}{Baden-W\"{u}rttemberg}{70569}{Germany}
\paperauthor{Christof Iserlohe}{christof.iserlohe@irs.uni-stuttgart.de}{0000-0003-4223-7439}{Universit\"{a}t Stuttgart}{Institut f\"{u}r Raumfahrtsysteme}{Stuttgart}{Baden-W\"{u}rttemberg}{70569}{Germany}
\paperauthor{Bastian Knieling}{bastian.knieling@irs.uni-stuttgart.de}{0009-0000-7038-0406}{Universit\"{a}t Stuttgart}{Institut f\"{u}r Raumfahrtsysteme}{Stuttgart}{Baden-W\"{u}rttemberg}{70569}{Germany}
\paperauthor{Mario Ruoff}{mario.ruoff@irs.uni-stuttgart.de}{0009-0006-6676-4865}{Universit\"{a}t Stuttgart}{Institut f\"{u}r Raumfahrtsysteme}{Stuttgart}{Baden-W\"{u}rttemberg}{70569}{Germany}
\paperauthor{Bernhard Schulz}{bernhard.schulz@irs.uni-stuttgart.de}{0009-0006-7941-9809}{Universit\"{a}t Stuttgart}{Institut f\"{u}r Raumfahrtsysteme}{Stuttgart}{Baden-W\"{u}rttemberg}{70569}{Germany}

% There should be one \aindex line (commented out) for each author. These are used to
% build up the author index for the Proceedings. The surname must come first, followed by
% initials. Note the use of ~ before each initial to control spacing.
% The \author entries (see above) have surname last. These \aindex entries have
% surname first.
% The Aindex.py command willl create them for you after you have constructed the \author
% The first entry should be the first author, for bold-facing the author index page-reference

%\aindex{Bryant,~A.}
%\aindex{Fr\"{u}h,~J.}
%\aindex{Greiner,~B.}
%\aindex{Iserlohe,~C.}
%\aindex{Knieling,~B.}
%\aindex{Ruoff,~M.}
%\aindex{Schulz,~B.}


\begin{abstract}
In 1991, Steve Lord released the ATRAN tool for computing atmospheric transmission
spectra from the near-infrared to microwave wavelengths. For decades, ATRAN has
served as a standard for telluric correction of ground-based and airborne astronomical
observations, including several of the infrared instruments that flew aboard the
now-decommissioned airborne observatory SOFIA.

At the SOFIA Data Center (SDC, see \citet{P604_adassxxxvi}), revised versions of
the original ATRAN code and its spectral line database have been developed, with
the permission of the author, incorporating upgrades and fixes matching the specific
needs of SOFIA science data. This new model, named ATRAN\_SDC, has been used to
generate a databank of mid- and far-infrared transmission spectra, with parameters
corresponding to the typical atmospheric conditions during SOFIA science operations.
These spectra are publicly available on the University of Stuttgart's Data Repository,
and are utilised by the data reduction pipeline package "Redux''.

To further improve the access of the global community to these valuable data assets,
and adhere to the ``FAIR'' data principles, the SDC has developed a standalone web application.
This interfaces with the ATRAN\_SDC code, and allows users to generate fully customised
ATRAN\_SDC transmission spectra, with a range of relevant parameters, and download them
for further analysis and usage on infrared datasets.

This article overviews the development of the ATRAN\_SDC transmission model and
transmission spectra, as well as the functionality of the associated web application.
It also explores the integration of the updated spectra into the Redux data
reduction pipeline, as part of an enhanced telluric correction for the mid- and
far-infrared instruments FORCAST and FIFI-LS.
\end{abstract}

% These lines show examples of subject index entries. At this stage these have to commented
% out, and need to be on separate lines. Eventually, they will be automatically uncommented
% and used to generate entries in the Subject Index at the end of the Proceedings volume.
% Don't leave these in! - replace them with ones relevant to your paper.

% These lines show examples of ASCL index entries. At this stage these have to commented
% out, and need to be on separate lines. Eventually, they will be automatically uncommented
% and used to generate entries in the ASCL Index at the end of the Proceedings volume.
% The ascl.py command will scan your paper on possible code names.
% Don't leave these in! - replace them with ones relevant to your paper.


\section{The ATRAN Transmission Model}
In the near- to far-infrared wavelength regimes, such as those covered by the Stratospheric
Observatory for Infrared Astronomy (SOFIA) during
its operation, the opacity of the Earth's atmosphere plays a key role in the feasibility
of ground-based and airborne observations. Among the numerous atmospheric
transmission models is ATRAN (Atmospheric TRANSmission), developed by Steve Lord (\citet{P809_Lord_1992}). ATRAN
computes the radiative transfer through a number of model layers characterised by
constant conditions and absorbant abundances. By specifying an observatory altitude,
zenith angle, water vapour overburden, and ozone latitude profile, a transmission spectrum
can be generated covering wavelengths from 0.8 \textmu m to 10 mm.

ATRAN was used for telluric correction on much SOFIA science data during its operation,
including that of the far-infrared spectrometer FIFI-LS (\citet{FIFI_Fischer_2021},
\citet{FIFI_Iserlohe_2021}, \citet{FIFI_Iserlohe_2022}), where it has been integrated
into the data reduction package Redux (see \citet{P816_adassxxxvi} for further details).
\citet{FIFI_Fadda_2023} compared transmission spectra derived from ATRAN with those of
another model, AM (\citet{P809_Paine_2024}), along with the observed telluric profile from FIFI-LS. In that work, several
wavelength regions were identified where ATRAN diverges from both the observed profile and
AM-modelled spectrum. This is shown in Figure \ref{P809_f1} below.

By applying the corresponding transmission spectrum to a science observation, the wavelength
regions of divergence will be significantly over-corrected. At the time, no obvious cause
for the divergence was found or postulated. However, the SDC has since identified that the 1998
release of the HITRAN molecular aborption database used by ATRAN exhibits a faulty scaling of
the molecular oxygen isotopologue $^{17}$O$^{16}$O abundance (see \citet{ROTHMAN1998665}),
which is responsible for the issues identified in \citet{FIFI_Fadda_2023}.

\articlefigure[width=0.8\textwidth]{P809_f1.png}{P809_f1}{Plot of observed FIFI-LS telluric profile (blue), and AM (orange) and
ATRAN (green) model transmission spectra. Regions around 139.3, 147.7 and 165 \textmu m demonstrate
divergence of ATRAN from the AM and observed spectra.}

\section{The ATRAN\_SDC Database}
With the permission of Steve Lord, the SDC has built an updated version of the ATRAN
model, incorporating a number of alterations and fixes, including for the aforementioned
isotopologue error. This model is known as ATRAN\_SDC, and has been used to regenerate
a full databank of transmission spectra covering mid- to far-infrared wavelengths,
parameterised by typical SOFIA flight altitudes, telescope zenith angles, and water vapour
overburdens.

While the underlying ATRAN codebase remains the property of Steve Lord, the databank of
ATRAN\_SDC spectra has been made available to the community via the Data Repository of the
University of Stuttgart (DaRUS)\footnote{https://darus.uni-stuttgart.de/dataverse/irs-sofia-ad/?q=ATRAN},
where they are also automatically retrieved by the latest Redux reduction software for the
telluric correction of FIFI-LS and FORCAST spectra. This is a significant improvement
to previous versions of the pipeline, which required users to download ATRAN spectra prior
to reduction. Local transmission spectra are still supported, but now telluric correction
using the latest ATRAN\_SDC data requires only the files stored on the SDC GitHub
repository\footnote{github.com/SOFIA-Data-Center/sofia\_redux}.

\section{The ATRAN\_SDC Web Application}
In order to further improve the accessibility of ATRAN\_SDC spectra to the community and
expand their application beyond the correction of SOFIA science data, the SDC has developed
a web-based application\footnote{https://atran.sofiadatacenter.de/} for the generation,
analysis and download of transmission spectra. This also restores the functionality of the
former web-based ATRAN server, which disappeared after the end of SOFIA operations.
The application has been tested for support of all modern desktop and mobile operating systems.
A Javascript interface allows for interactive manipulation of all relevant model and display parameters.
The Django-based backend queries the ATRAN\_SDC model, and displays the resulting spectrum in an interactive
plot window, as shown in Figure \ref{P809_f2} below.

Multiple spectra with different parameters can be overplotted for
comparison, along with other useful display options such as spectral emission line locations.
Molecular absorbants can be displayed individually, and exported as .svg plots.
Spectra can be downloaded as .FITS files, compatible with the Redux reduction
suite, or as .dat tables. Parameters also allow customisation of the water vapour and
ozone overburden values, extending the model to conditions expected in some of the
higher ground-based observatories. This may be of particular interest to the mid- and far-infrared
communities, opening up access to these wavelengths in specific windows under optimal weather.

\articlefigure[width=0.89\textwidth]{P809_f2.png}{P809_f2}{Screenshot of the ATRAN\_SDC Web Application, with example transmission spectra displayed for two different parameter sets.}

\clearpage

\acknowledgements The SDC is funded by the DLR Space Administration under grant number FKZ 50OK2404.
The Authors are grateful to the creator and maintainer of ATRAN, Steve Lord, for his support and
permission to use his program in the manner presented here.

\bibliography{P809}  % For BibTex

% if we have space left, we might add a conference photograph here. Leave commented for now.
% \bookpartphoto[width=1.0\textwidth]{foobar.eps}{FooBar Photo (Photo: Any Photographer)}

\end{document}
